% Generated draft: requires mathematical and source review.
\documentclass[11pt]{article}
\usepackage[margin=25mm]{geometry}
\usepackage{fontspec}
\setmainfont{Latin Modern Roman}
\usepackage{amsmath,amssymb,mathtools,unicode-math}
\setmathfont{Latin Modern Math}
\usepackage{enumitem,needspace,xurl,hyperref}
\setlength{\parindent}{0pt}
\setlength{\parskip}{5pt}
\setlength{\emergencystretch}{3em}
\newcommand{\N}{\mathbb{N}}
\newcommand{\Z}{\mathbb{Z}}
\newcommand{\Q}{\mathbb{Q}}
\newcommand{\R}{\mathbb{R}}
\newcommand{\C}{\mathbb{C}}
\newcommand{\F}{\mathbb{F}}
\newcommand{\A}{\mathbb{A}}
\newcommand{\PP}{\mathbb{P}}
\newcommand{\E}{\mathbb{E}}
\newcommand{\eps}{\varepsilon}
\newcommand{\dd}{\,\mathrm{d}}
\begin{document}
\section*{UnsolvedMath ID: 800169}
\textbf{Source problem code:} catalogue entry
\subsection{Definitions and mathematical statement}
Under the usual interpretation of primitive forms here as holomorphic cuspidal Hecke eigenforms, let \(\mathcal H_k\) denote the set of primitive forms of level \(1\) and even integral weight \(k\geq 2\), with each form normalized to have first Fourier coefficient \(1\). For \(f\in\mathcal H_k\) and \(g\in\mathcal H_\ell\), let \(L(s,f\otimes g)\) denote their Rankin--Selberg \(L\)-function in the normalization with central point \(s=\tfrac12\).

Fix a constant \(C\geq0\), independent of \(K\). A natural unweighted interpretation of the sixth moment in the supplied question is
\[
M_6(K;C)=
\sum_{\substack{k,\ell\in 2\Z_{\geq1}\\ K<k+\ell\leq2K\\ |k-\ell|\leq C}}
\ \sum_{f\in\mathcal H_k}\ \sum_{g\in\mathcal H_\ell}
\left|L\!\left(\tfrac12,f\otimes g\right)\right|^6.
\]
The fixed constant \(C\) makes explicit the source condition \(|k-\ell|=O(1)\).

\textbf{Requested task.} Use the recent work of Humphries--Khan, cited as [HK] in the supplied source, to obtain upper bounds in the weight aspect for the sixth moment of these central values as \(K\to\infty\), subject to the displayed restrictions on \(k\) and \(\ell\).

\textbf{Unresolved qualifications.} The displayed sum is an explicit possible interpretation, not a uniquely specified moment in the source. The record does not state whether the moment is unweighted or harmonically weighted, whether it is a sum or a normalized average, or whether the bound is intended for each admissible weight pair or after summation over weights. It gives no target exponent or precise dependence on the bounded weight difference. The holomorphic cuspidal interpretation and the normalization above make the notation explicit but require confirmation against the original source. The bibliographic identity and applicable hypotheses of [HK] are not supplied.
\subsection{Short English statement}
Alia Hamieh asks for upper bounds for sixth moments of central Rankin--Selberg values attached to level \(1\) primitive forms whose weights differ by a bounded amount and satisfy \(K<k+\ell\leq2K\), using work of Humphries--Khan. The averaging convention and the precise reference [HK] remain unresolved. The supplied registry labels the problem open; this is not an independently verified status assessment.
\subsection{Sources}
\begin{itemize}
\item Local UnsolvedMath registry: permanent ID \(800169\), source problem code \texttt{catalogue entry}. The supplied record labels the problem \emph{open} and records Hamieh, Kadiri, Ng as proposers, with proposed year \(2024\).
\item Hamieh, Kadiri, Ng, \emph{Moments of L-functions Problem List} (2024), Problem 19. Supplied source URL: \url{https://arxiv.org/abs/2405.17800}. The record reports access on 2026-10-16 and extraction by \texttt{arxiv\_latex}; these are supplied provenance details, not a claim of independent access or verification.
\item The problem statement specifically attributes the question to Alia Hamieh, UNBC.
\item Humphries--Khan [HK]: the supplied statement refers to their recent work, but supplies no title, publication details, URL, or theorem reference identifying that work.
\end{itemize}
\end{document}
